\pdfoutput=1
\documentclass[11pt]{article}
\usepackage[margin=1in]{geometry}
\usepackage{amsmath,amssymb,amsthm}
\usepackage{enumitem}
\usepackage{booktabs}
\usepackage{array}
\usepackage{longtable}
\usepackage{hyperref}
\hypersetup{hidelinks,
  pdftitle={Questions on the remaining value-group profile switch from Theorem 3.11 to Corollary 3.12: where is it supplied? A central question with a test, and component questions on local shears and object membership},
  pdfauthor={Toshihisa Urahigashi}}

\newcommand{\uu}[1]{\underline{\underline{#1}}}
\newcommand{\ord}{\operatorname{ord}}
\newtheorem{question}{Question}
\theoremstyle{definition}
\newtheorem{statement}{Statement}
\newtheorem{remark}{Remark}
\newtheorem*{centralq}{Central question}

\title{Questions on the remaining value-group profile switch\\
from Theorem~3.11 to Corollary~3.12\\[4pt]
\large Where is it supplied? A central question with a test,\\
and component questions on local shears and object membership}
\author{Toshihisa Urahigashi\\
\small\texttt{123026.urahigashi.toshihisa@gmail.com}}
\date{Version 2.0.0 --- 14 November 2026\\[3pt]
\small Independent researcher. This work was carried out personally and is
unrelated to any institutional affiliation.\\[3pt]
\small ORCID \href{https://orcid.org/0009-0004-0460-6242}{\texttt{0009-0004-0460-6242}}\\[-2pt]
\small \href{https://doi.org/10.5281/zenodo.24213391}{doi:10.5281/zenodo.24213391} (all versions: \href{https://doi.org/10.5281/zenodo.24198744}{10.5281/zenodo.24198744})}

\begin{document}
\maketitle

\begin{abstract}
In Step~(xi) of the proof of Corollary~3.12 of inter-universal
Teichm\"uller theory, the $q$-pilot is compared with the possible
images of the $\Theta$-pilot produced by the multiradial algorithm
of Theorem~3.11, and the inequality is drawn from the property
(SHE) in the form stated in Step~(xi-e).  Several works, in
abstractions of the passage and, in the author's earlier notes, for
the statements that the proof actually cites, show that this form
does not follow from what is stated about the output of the
algorithm; in a conditional assembly, the remaining input is a
value-group profile switch whose profile part is independent of
every other input.  Whether the construction of the algorithm
itself supplies the step, and if so where, is not decided by those
results.  This version of the note asks that question directly.

Together with IUT~IV, the step has the strength of an explicit
height inequality.  Every proof of such an inequality has a step of
this strength, at the place where its content enters.  The fact is
not offered as evidence for or against the existence of the
construction; it is used only as a criterion by which a supplier,
once found, can be recognized.  The question asks where, not how,
and it can be answered by pointing to a statement or to a passage
of a proof.

The step is the value-group counterpart of a compatibility that is
proved on the unit side, where cyclotomic rigidity gives a unique
isomorphism.  On the value-group side the two profiles differ by the
factors $j^2$ between the labels, and the counterpart cannot be an
isomorphism: in that form it would force the degree scale of the
$\Theta$-link to be $1$.  It has to be an inclusion in the
holomorphic hull, and that inclusion is the inequality.
Section~\ref{s:central} locates where the source divides the two
sides.  Step~(i) of the proof splits the $\Theta$-link into a unit
part and a value part, and the compatibilities of Theorem~3.11(iii)
concern the unit part; we found no statement that carries the value
part to regions in the common container.  Since the proof of
Theorem~3.11 states that its assertions follow immediately from the
definitions and the cited references, the question becomes which
definition or cited construction supplies that one edge.  We read
five places that could supply it, the family of possible images,
the holomorphic hull, the identification of theta and Gaussian
monoids, the compatibility statement of Theorem~3.11(iii)(c) and the
global realified Frobenioids of Step~(xii), as holding in the
author's models of the cited statements, built from actual local
objects, in which the step fails.  This is a reading, not a claim
that the construction does not supply the step.  The question is
which property of the actual construction excludes those models,
that is, where the models fail to represent it, or which of the
readings is wrong.  The questions are not addressed to any particular person
or project.

The small questions Q0--Q7 of earlier versions, on local shears,
their globalization and object membership, are kept in appendices
as component questions about one route, through the difference of
the theta and algebraic sections of [EtTh].  This version answers
them provisionally: the difference of sections does not move degree
readouts, a degree-moving freedom is one positive scalar common to
all places, and the route does not carry the profile switch.  The
abstract implication cores used by the component questions are
checked in Lean, as before, and a standalone Lean file checks the
elementary forms that the readings use; the readings themselves are
not formalized.

No claim is made that Corollary~3.12, the ABC conjecture or the
profile switch is false, or that the construction does not supply
the step.

\medskip
\noindent\textbf{Keywords:} inter-universal Teichm\"uller theory,
Corollary~3.12, Step~(xi), simultaneous holomorphic expressibility,
value-group profile, local shear, mono-theta environment,
anabelian geometry, Lean.
\end{abstract}

\medskip
\noindent\textbf{Note added in version 2.0.0.}
This version adds Section~\ref{s:central}, the central question, and
makes it the main question of the note; Section~\ref{ss:square}
locates the one edge of the comparison for which no statement was
found, Section~\ref{ss:test} describes the author's models, reads
five places that could supply the edge as holding in them, and asks
which property of the actual construction excludes them or which
reading is wrong, and Appendix~\ref{s:questions} gives provisional answers to
Q0--Q7; a standalone Lean file checks the elementary forms that the
readings use (Section~\ref{ss:cores}).  The abstract, the position
statement and the closing section are rewritten accordingly; each
component question receives a line on its bearing on the central
question; the sentence of earlier versions saying that none of the
questions is Corollary~3.12 or ABC is replaced by
Remark~\ref{r:strength}; earlier works on the same step are cited
([Higuchi], [Pajo], [Takkun], [Saito], [sfingali], [Geiger], [Form],
[Ri]); and the references to the author's notes are updated.  The
component questions Q0--Q7, Statement~P$'$ and the related tables
are moved to Appendices~\ref{s:setup}--\ref{s:relations}; apart from
the lines on their bearing, a paragraph on their scope and the
provisional answers, their content is unchanged.  Version~1.0.2 rewrote the comments of the Lean files in
English and changed no question.

\section*{Position of this note}

This note is a problem statement for joint examination, not a
completed proof.  Its questions are not addressed to any particular
individual or project; they are addressed to anyone with relevant
expertise who can supply a reference, a correction, a
counterexample or a better formulation.  Answers can be sent to the
e-mail address on the title page or posted as an issue in the
public repository that holds this note,
\url{https://github.com/turahigashi/iut-explicit-initial-theta}; a one-line answer in one of
the forms of Section~\ref{ss:answers} is enough.  The note does not
claim a proof or disproof of Corollary~3.12 or the ABC conjecture, a
determination of H*val, or an equivalence between Statement~P$'$
below and the residual.  It is not an alternative to, or a
criticism of, Project LANA's (9-1): (9-1) is one of the reference
forms of the step asked about in Section~\ref{s:central}, and the
component questions separate definitions, local mechanisms and
globalization steps that precede or surround it.  Even if every
component question receives a positive answer, the two return
bridges and the numbered IUT inputs remain separate.  Conversely, a
concrete countermodel to one component question would locate a
place where this route closes.

\section{The central question}\label{s:central}

\subsection{The step}\label{ss:step}

Corollary~3.12 of [IUTchIII] is proved in twelve steps, and the
inequality is drawn in Step~(xi) [IUTchIII, pp.~181--185].  We
summarize the substeps that matter here.  The summary is our
reading and is meant to be checked.
\begin{itemize}[leftmargin=4.2em,labelwidth=3.6em,labelsep=0.6em,align=left]
\item[(xi-a)] The $\Theta_{\mathrm{LGP}}^{\times\boldsymbol\mu}$-link
from $(0,0)$ to $(1,0)$ is read as a gluing isomorphism under which
the $\Theta$-pilot object at $(0,0)$, regarded as an object of the
relevant global realified Frobenioid, corresponds to the $q$-pilot
object at $(1,0)$.
\item[(xi-b)] The multiradial algorithm of [IUTchIII, Thm.~3.11]
constructs a collection of possibilities for output data in a
container common to the two columns, with the properties (IPL) and
(SHE) of [IUTchIII, Rem.~3.11.1(iii)].  The author states that at
this step only the qualitative properties (IPL), (SHE) and (APT) of
the algorithm will be used, ``temporarily forgetting'' that it
involves the $\Theta$-pilot, theta functions and mono-theta
environments, which form its hidden internal structure in the sense
of [IUTchIII, Rem.~3.11.1(iv)].
\item[(xi-c)] The possibilities for output data are related,
through isomorphisms of
$\mathcal F^{\Vdash\blacktriangleright}$-prime-strips, to the
representation of the $q$-pilot obtained by the column-one
log-Kummer correspondence (IPL), and their construction is
expressible in terms of the holomorphic structure of column one
alone (SHE).  After passing to the holomorphic hull they are
represented by localizations of arithmetic vector bundles.
\item[(xi-d)] Taking determinants of suitable positive tensor
powers and procession-normalized log-volumes yields the half-line
$\mathbb R_{\le-|\log(\uu\Theta)|}$ of possible values and the
number $-|\log(\uu q)|$ attached to the $q$-pilot.  The column-one
log-Kummer correspondence is used here to correct the vertical
misalignment of the unit parts.
\item[(xi-e)] (SHE), that is, expressibility in terms of the
arithmetic holomorphic structure of column one, is read as
expressibility by operations that remain valid when the input
prime-strip is subject to the condition that its pilot log-volume
equals the fixed value $-|\log(\uu q)|$.
\item[(xi-f)] Hence the construction of the possible values is also,
perhaps only up to a kind of approximation caused by the
indeterminacies, a construction of $-|\log(\uu q)|$, so that
$-|\log(\uu q)|\in\mathbb R_{\le-|\log(\uu\Theta)|}$.  This is the
inequality.
\item[(xi-g)] The argument is summarized as two tautologically
equivalent ways of computing the log-volume of the $q$-pilot
[IUTchIII, Fig.~3.8].
\end{itemize}
The question of this note concerns the passage from (xi-e) to
(xi-f): what supplies (SHE) in the form used in (xi-e).

\subsection{What is already known}\label{ss:known}

\begin{enumerate}[label=(K\arabic*),leftmargin=*]
\item \emph{The step against the cited statements.}  At the level of
log-volumes, the form of (SHE) in (xi-e) is the conclusion
rewritten; it is equivalent to the inequality [U26, Prop.~6].  Its
weak reading, that the possible images are computed from the common
container, follows from the statements that the proof cites, taken
at the strength of the source [U26, Prop.~5].  The strong readings,
and the equality (9-1) that the interim report of Project LANA
identifies as the missing step [LANA, \S9.2], do not
[U26, Thm.~10, Cor.~11].  One of the models that show this is built
from actual local objects at a finite extension of $\mathbb Q_p$ carrying
the local data of a theta function: tensor packets of the field
with the mono-analytic integral structure, the indeterminacies
(Ind1) and (Ind2) acting through all permutations of tensor factors
and all automorphisms of the log-shells, values of the theta
series, a container containing all horizontal and vertical images,
and, as hull, the smallest multiple of the ring of integers of the
packet that contains the container [U26, Thm.~22].  The models
survive the addition of any finite set of further places [U26, Thm.~14],
and the closed loop of [IUTchIII, Rem.~3.9.5(viii)(ix)], to which
[IUTchIII, Rem.~3.11.1(iii)] relates (SHE), holds in them.
\item \emph{The step in a conditional assembly.}  In [U26c] the chart
calibration, the normalized-generator scale and the unit side are
supplied at the level of stated inputs, and the remaining input is
the value-group profile switch H*val: from the
$\{j^2u\}_j$-type profile of the $\Theta$-pilot, an actual $q$-pilot
representation compatible with the holomorphic hull is obtained
[U26c, \S8].  The model of [U26d], with the source profile at
$\ell=7$, satisfies every other input of the assembly and falsifies
the conclusion.  In the language of [U26b], the step is a failure
of completeness at an interface, and the question becomes whether
the additional principle needed is supplied by the actual
construction.
\item \emph{Earlier work on the same passage.}  Higuchi [Higuchi]
proves non-derivability for a qualitative package extracted from
(xi-e) and (xi-f), and Pajo [Pajo] a no-go theorem for a scale-free
monoid, hull and log-volume abstraction.  In Lean, Takkun-kohinata
[Takkun] proves for its transcription of the passage that the
volume bound of the multiradial representation together with the
containment is equivalent to the height inequality, that strict
identification of the copies, with volume-preserving
indeterminacies and a hull no larger than the $\Theta$-pilot, is
contradictory, and that the slack between the hull and the
$\Theta$-pilot cannot come from the orbits of (Ind1) and (Ind2),
only from the hull.  Saito [Saito] audits the passage from
(xi-a)--(xi-e) to the membership of (xi-f) against a registered
inventory of the source; sfingali [sfingali] formalizes in Lean the
core of the objection of [SS]; and Geiger [Geiger] records that the
decisive membership of (xi-f) is stated to follow formally from a
pair of declared properties rather than derived.  The public Lean
repository of Project LANA [lana] proves that Step~(xi) cannot be
derived from a package collecting what the repository proves at the
level of Theorem~3.11, shows that Step~(xi) adds to the package
exactly the inequality $(s-1)a\le V$ in its notation
(\texttt{xi\_iff}), and separates two readings of (SHE), of which
the one asserting that some output possibility has the volume class
of the $q$-pilot region is proved equivalent to the wall inequality
of [LANA, \S0.2] (\texttt{sheq\_iff\_wall}).  The interim report of the project
states that the project does not at present have a proof of (9-1)
[LANA, \S10.5].
\item \emph{The author of IUT on the same step.}  In his report on
formalization [Form, \S4], Mochizuki reorganizes the implication
from Theorem~3.11 to Corollary~3.12 into four commutative triangles
and isolates an intermediate statement ``3.11.5'' (Theorem~3.11
together with Remark~3.9.5), which, as the report notes, does not
appear in [IUTchI]--[IUTchIV].  The skeletal Lean code described
there for ``3.11.5 $\Rightarrow$ 3.12'' takes as an input an arrow
(SHE), which takes the input data of the algorithm to be the data of
the $q$-pilot constructed from the \'etale Hodge theatre.  The report
also states that the implication from Theorem~3.11 to
Corollary~3.12, which has attracted the most attention, is in fact
not very difficult, and that, correctly understood, it shows that the
real object of interest is the proof of Theorem~3.11 [Form, \S2].  At the level of log-volumes
the arrow (SHE) is the carrier reading of [U26, Def.~4]
[U26, \S10].  Richards' Lean port [Ri] states ``3.11.5'' as a
proposition and derives it, in a separate file, from a structure
with a field identifying the region of the $\Theta$-pilot with that
of the $q$-pilot.
\end{enumerate}

\subsection{What the step must deliver}\label{ss:target}

We state it in terms of local fields, log-shells and Haar measure.
Fix a number field $K$, an elliptic curve $E/K$, a prime $\ell\ge5$
and the remaining data of an initial $\Theta$-datum
[IUTchI, Def.~3.1]; further notation is that of
Appendix~\ref{s:setup}.  At a bad place $v$ put
$u_v=|\log(\uu q_v)|=(2\ell)^{-1}\log(\mathfrak q_v)$
[IUTchIV, Thm.~1.10].  The labels are
$j=1,\ldots,\ell^\divideontimes=(\ell-1)/2$.  At $v$ the
$\Theta$-pilot has the labelled log-volume profile
$\{-j^2u_v\}_j$ and the $q$-pilot the profile $\{-u_v\}_j$
[IUTchII, Cor.~4.5(iv)(v)], [IUTchIII, Def.~3.8(i)].  The procession
normalization gives every capsule the same weight
[IUTchIII, Prop.~3.9(i), Rem.~3.9.3], and the aggregates of the two
profiles are $-\frac{\ell(\ell+1)}{12}u_v$ and $-u_v$
[U26d, \S3].

The possible images of the $\Theta$-pilot are its images under the
indeterminacies of [IUTchIII, Thm.~3.11]: (Ind1), coming from
automorphisms of the \'etale-like data (copies of the local absolute
Galois groups), which act on log-shells through mono-anabelian
reconstruction; (Ind2), coming from automorphisms of the unit parts;
and (Ind3), the upper semi-compatibility of the log-Kummer
correspondence [IUTchIII, Thm.~3.11(ii)].  The holomorphic hull of
their union is formed in tensor packets of log-shells, and volumes
are procession-normalized log-volumes with respect to Haar measure
[IUTchIII, Rem.~3.9.5(vii)].  Explicit formulas for (Ind1)--(Ind3)
in terms of local fields are given in [DHPS].

The step must deliver the following.
\begin{itemize}[leftmargin=2.6em,labelwidth=2em,labelsep=0.6em,align=left]
\item[(S)] The region of the $q$-pilot, represented in the common
container through the column-one log-Kummer correspondence, lies in
the holomorphic hull of the possible images of the $\Theta$-pilot
(\emph{carrier form}); or at least its log-volume does not exceed
that of the hull (\emph{value form}):
\[
-|\log(\uu q)|\ \le\ -|\log(\uu\Theta)|.
\]
\end{itemize}
The value form is the inequality of Corollary~3.12.  In the reading
of [U26, Def.~4], the carrier form implies (9-1), and (9-1) implies
the value form [U26, Prop.~6].  Together with the upper bound for
$-|\log(\uu\Theta)|$ in [IUTchIV, Thm.~1.10], (S) for all initial
$\Theta$-data gives, through [IUTchIV, \S2], the diophantine
inequalities of [IUTchIV, Cor.~2.3].

\subsection{What is known not to suffice}\label{ss:notsuffice}

\begin{enumerate}[label=(N\arabic*),leftmargin=*]
\item \emph{The log-volume content of the cited statements.}  The
22 statements that the proof of Corollary~3.12 uses as definitions
or premises, among 63 cited targets [U26, \S3, App.~A], include what
[IUTchIII, Thm.~3.11] states about the output of the algorithm, the
hull of [IUTchIII, Rem.~3.9.5], (IPL) and the closed loop.  Their
content at the level of log-volumes holds in the models of (K1),
described in Section~\ref{ss:test}, in which (S) fails.  So does the weak reading of (SHE).
\item \emph{Scalar calibrations.}  Agreement of the measuring charts
of the two routes [U26c, Lemma~8.1], agreement at the normalized
generator $j=1$, and compatibility on the unit side
[U26c, Lemma~8.2] hold in the model of [U26d], in which (S) fails.
A separate assumption that the aggregate ratio of the two pilots
equals $1$ contradicts the source readout $\ell(\ell+1)/12\ne1$
[U26d, \S7.5].
\item \emph{Identification of the two pilots.}  A statement
identifying the image of the $\Theta$-pilot with the region of the
$q$-pilot implies (S), but it forces the degree scale of the
$\Theta$-link to be $1$ [U26, Prop.~26(b)], contrary to the source
readout; with the source data it cannot hold.  Compare [SS],
[sfingali] and \texttt{strict\_collapse} in [Takkun].
\end{enumerate}
In the language of [SL], a statement of type (N3) asks for a strict
2-cell between the two paths, which exists only if the degree scale
is $1$ [SL, Thm.~5.2]; (S) is a lax 2-cell [SL, Thm.~5.3]; and the
properties (IPL), (SHE) and (APT) by themselves are compatible with
any size of the blurring, so that neither (S) nor its negation
follows from them [SL, Thm.~5.7].

\subsection{The value-group counterpart, and the edge without a statement}\label{ss:square}

The step is the value-group counterpart of a compatibility that is
proved on the unit side.  On the unit side the indeterminacies (Ind1)
and (Ind2) preserve log-volumes [IUTchIII, Cor.~3.12, proof,
Step~(x)], and what is needed is an isomorphism, made unique by
cyclotomic rigidity [EtTh, Cor.~2.19(i)].  In the conditional
assembly this side is supplied by such uniqueness statements
[U26c, Lemma~8.2], and [EtTh, Cor.~2.19(i)] is also the solved form
from which Statement~P$'$ departs (Appendix~\ref{s:pprime}).  On the
value-group side the two profiles differ by the factors $j^2$ between
the labels.  A counterpart of the same type, an isomorphism
compatible with log-volumes, would identify the two pilots and force
the degree scale of the $\Theta$-link to be $1$, as in (N3).  What is
needed is an inclusion in the hull, (S): a lax rather than a strict
2-cell [SL, Thms.~5.2, 5.3].  This is why the step does not follow by
analogy from the unit side, and why it carries the strength of the
conclusion.

In the column-one container the proof compares what is carried along
two routes.
\begin{itemize}[leftmargin=4.8em,labelwidth=4.2em,labelsep=0.6em,align=left]
\item[Route 1] The $\Theta$-pilot at $(0,0)$, the value part of the
$\Theta^{\times\boldsymbol\mu}_{\mathrm{LGP}}$-link, the $q$-pilot at
$(1,0)$, and the Kummer isomorphisms of column one
[IUTchIII, Thm.~3.11(ii)].  The statement of Corollary~3.12 defines
$-|\log(\uu q)|$ as the log-volume of this image, computed in the
same multiradial representation and not subject to (Ind1)--(Ind3)
[IUTchIII, p.~174].
\item[Route 2] The $\Theta$-pilot at $(0,0)$, the Kummer isomorphisms
of column zero, and the permutation symmetry at the end of
[IUTchIII, Thm.~3.11(i)], which carries the possible images into the
column-one container (Step~(xi-b)).
\end{itemize}
Step~(i) of the proof splits the $\Theta$-link into a unit part, a
full poly-isomorphism of $\mathcal F^{\vdash\times\boldsymbol\mu}$-prime-strips,
and a value part, a full poly-isomorphism of
$\mathcal F^{\Vdash\blacktriangleright}$-prime-strips that sends the
$\Theta$-pilot to the $q$-pilot [IUTchIII, p.~175].  The edges of the
comparison and the statements that supply them are as follows.

\begin{center}
\small
\begin{tabular}{@{}>{\raggedright\arraybackslash}p{0.50\textwidth}>{\raggedright\arraybackslash}p{0.44\textwidth}@{}}
\toprule
Edge & Statement\\
\midrule
Kummer isomorphisms of both columns & [IUTchIII, Thm.~3.11(ii)]\\
Permutation symmetry between the columns & [IUTchIII, Thm.~3.11(i)], end\\
Compatibility with the unit part of the $\Theta$-link & [IUTchIII, Thm.~3.11(iii)(a)(b)]\\
Shielding of the cyclotomes that give the value part & [IUTchIII, Thm.~3.11(iii)(c)(d)]; Step~(ii) of the proof\\
Log-volumes, their compatibility with log-Kummer, the hull & [IUTchIII, Thm.~3.11(ii), end; Prop.~3.9(iv); Rem.~3.9.5]\\
Value part of the $\Theta$-link, carried to regions in the container & none located\\
\bottomrule
\end{tabular}
\end{center}

The compatibilities of [IUTchIII, Thm.~3.11(iii)] are stated with the
full poly-isomorphisms of $\mathcal F^{\vdash\times\boldsymbol\mu}$-prime-strips
as the example, in the sense that the constructions are stable,
equivariant and functorial with respect to arbitrary automorphisms of
the domain and the codomain of the $\Theta$-links
[IUTchIII, pp.~157--158].  Step~(xi-b) describes the relation between
the value parts as consisting precisely of full poly-isomorphisms of
$\mathcal F^{\Vdash\blacktriangleright}$-prime-strips, and Step~(xi-e)
crosses the remaining edge by its reading of (SHE).  The proof of
[IUTchIII, Thm.~3.11] states that each assertion follows immediately
from the definitions and the references cited in it
[IUTchIII, p.~159].  So if something supplies the edge, it lies in
those definitions and cited constructions, or in the hidden internal
structure that Step~(xi-b) sets aside.  This is our reading of the
statements, and the question below asks for it to be confirmed or
corrected.

\subsection{The models, the test, and the question}\label{ss:test}

\paragraph{The models.}  The models used in this section are the
author's own, from [U26] and [U26d].  They are countermodels to a
derivation: the statements in question hold in them, and (S) fails.
They are not counterexamples to Corollary~3.12, and they are not
actual global initial $\Theta$-data.
\begin{enumerate}[label=(M\arabic*),leftmargin=*]
\item \emph{The local model of (K1)} [U26, Def.~20, Thm.~22].  Let
$L$ be a finite extension of $\mathbb Q_p$ that carries a theta
datum: $\tau,i\in L^\times$ with $i^2=-1$ and $|\tau|<1$, elements
$\xi_1,\xi_2,\xi_3$ with
$\xi_j^{\,7}=\vartheta(i)/\vartheta(\pm\tau^j i)$ for the theta
series $\vartheta$ of [EtTh, Prop.~1.4], and $r$ with
$r^{14}=\tau^2$.  The Tate parameter is $q=\tau^2$, so that
$\xi_j^{\,14}=q^{j^2}$ ($\ell=7$, labels $j=1,2,3$); for every
$N\ge1$ the $N$-th powers again form a theta datum, with Tate
parameter $q^N$.  At the stage $j$ the packet is the tensor product
$P_j$ of $j+1$ copies of $L$ over $\mathbb Q_p$, with ring of
integers $\widetilde{\mathcal O}$ and log-volumes normalized as in
[IUTchIV, Prop.~1.4(i)].  The $\Theta$-pilot is
$\xi_j^N\cdot\widetilde{\mathcal O}$ and the $q$-pilot is
$r^N\cdot\widetilde{\mathcal O}$.  The possible images are all
regions between a horizontal image and the container.  The
horizontal images are the images of the $\Theta$-pilot under all
$\mathbb Q_p$-linear automorphisms of $P_j$ that preserve the
log-shell, all unit multiples and all theta data with the same
$\tau,i$; the container is the $p$-power multiple of the log-shell
packet that contains all horizontal and vertical images; and the
hull is the smallest multiple of $\widetilde{\mathcal O}$ that
contains the container.  The excess of the hull over the
$\Theta$-pilot is at most a constant $E_L$ that depends only on $L$,
whereas (S) requires $\frac{11}{42}N\,|\log|q||\le E_L$; so (S)
fails for large $N$.  The model satisfies the 22 cited statements
that carry log-volume content, and survives the addition of any
finite set of further places [U26, Thms.~10, 14, 22].  Every choice
in it that the source does not dictate makes the hull larger, and so
works against the failure of (S) [U26, Rem.~23].
\item \emph{The profile countermodel} [U26d].  At $\ell=7$, with
labels $j=1,2,3$ at one bad place and equal weights $1/3$, the
$\Theta$-side readouts are $-j^2$ and the $q$-side readouts are
$-1$, as in the source.  The normalized generator, the unit side,
membership, monotonicity, the hull domain and both volume
decompositions hold, while the aggregates are $-14/3$ and $-1$, and
the conclusion of the shape of Corollary~3.12 fails.
\end{enumerate}
Both are given in finite, machine-checked form [U26, \S11],
[U26d, \S6].  The question of this section is, in other words, where
these models fail to represent the actual construction.

\paragraph{The test.}  Call a statement $P$, that is, a numbered statement or a passage of
a proof, in [IUTchI]--[IUTchIV] or in the papers they cite, a
\emph{candidate supplier} if it passes the following tests.
\begin{enumerate}[label=(T\arabic*),leftmargin=*]
\item \emph{It does work.}  Some consequence of $P$ for regions,
log-volumes and labels fails in (M1) or in (M2).  Otherwise $P$ falls under (N1) or (N2).
\item \emph{It keeps the profile.}  $P$ is consistent with the source
readout: it does not force the aggregate ratio $\ell(\ell+1)/12$ of
the two pilots to equal $1$.  Otherwise $P$ falls under (N3).
\item \emph{It is not the conclusion.}  $P$ is established without
using Step~(xi-f) or Corollary~3.12.
\end{enumerate}
Because (M1) is built from actual local objects and survives the
addition of places, a statement passing (T1) must use
something that the cited statements do not state at the level of
log-volumes: for instance a property of the hull or of the family of
possible images beyond those stated [U26d, \S7.6], global
information about the initial $\Theta$-datum, or the structure that
Step~(xi-b) sets aside.  A candidate that holds in (M1) and (M2) is
not a supplier, whatever else it does.  The tests are necessary
conditions, not sufficient ones.

\begin{centralq}
Five places in the source could supply the edge of
Section~\ref{ss:square}.  We read each of them as holding in (M1)
and (M2), in which (S) fails.  This is our reading of the
statements, not a claim that the construction does not supply the
step.
\begin{enumerate}[label=(\alph*),leftmargin=*]
\item \emph{The family of possible images}
[IUTchIII, Thm.~3.11(i); Props.~3.1--3.10].  (M1) reads it
maximally, as all regions between a horizontal image and the
container [U26, Def.~20]; the actual family lies in it.
\item \emph{The holomorphic hull} [IUTchIII, Rem.~3.9.5(vii)--(x)].
The hull of (M1) is the hull of the whole container; the actual
hull, the hull of the possible images [IUTchIII, Rem.~3.9.5(i)], lies
in it.  Even so, its excess over the $\Theta$-pilot is at most
$E_L$, while the excess that (S) requires grows linearly with $N$
[U26, Thm.~22(a)].
\item \emph{The identification of theta and Gaussian monoids}
[IUTchII, Cors.~4.5(iv)(v), 4.6(iv)(v)].  Its evaluation
isomorphisms, and the ratio vector $(\dots,j^2,\dots)$, are
constructed inside one Hodge theatre on the $\Theta$-side, and the
$q$-pilot does not appear in them; the $\Theta$-pilot of (M1) is
already built from the Gaussian values.
\item \emph{The compatibility statement}
[IUTchIII, Thm.~3.11(iii)(c)].  Its data are those of
[IUTchIII, Cor.~2.3]: a Hodge theatre, the
$\mathcal F^{\Vdash}$-prime-strip of the environment, the data at
the bad places, the unit part
$\mathcal F^{\vdash\times\boldsymbol\mu}_{\triangle}$ and a full
poly-isomorphism of unit parts.  The multiradiality of
[IUTchIII, Cor.~2.3(i)] is that the functor from these data to the
coric data
$(\mathcal D^{\vdash},\mathcal F^{\vdash\times\boldsymbol\mu}(\mathcal D^{\vdash}))$
is full and essentially surjective [IUTchIII, pp.~72--75].  The value
part $\mathcal F^{\Vdash\blacktriangleright}_{\triangle}$, which
carries the $q$-pilot, is in neither.  At the level of regions its
content is the closure under (Ind1) and (Ind2) and the bound (Ind3)
[U26, \S9], which hold in (M1); and it does not involve the value of
the $q$-pilot.
\item \emph{The global realified Frobenioids}
[IUTchIII, Cor.~3.12, proof, Step~(xii), p.~185].  Their role, as
stated there, is to fix whether a local object corresponds to a
given open submodule of the log-shell or to a $p^N$-multiple of it,
since the isomorphisms of Frobenioids that arise from the
$\Theta$-link preserve only isomorphism classes of objects
[IUTchIII, Rem.~3.6.2(i)].  In (M1) the regions and the
container are given subsets of the packet, at a fixed scale, so
this ambiguity does not arise.  The freedom that remains in the
degree is one positive scalar common to all places (Q0, Q1), which
acts on all labels alike and so cannot carry one profile to the
other.
\end{enumerate}
These readings depend on the container being a $p$-power multiple of
the log-shell packet [U26, \S6].  If they are right, (M1) and (M2)
satisfy (a)--(e).  Which property of the actual construction
excludes them, and where is it stated or constructed?  In other
words, where do these models fail to represent the actual
construction?  Or which of the readings (a)--(e) is wrong?

Our readings are least certain at the following places, and a
pointer on any one of them would be enough.
\begin{enumerate}[label=(\roman*),leftmargin=*]
\item What is the internal structure of the value parts that
Step~(xi-b) sets aside, and does it constrain regions?
\item Is the container a $p$-power multiple of the log-shell packet,
or is it larger, for instance growing with $N$?
\item Does the family of possible images constructed in
[IUTchIII, Props.~3.1--3.10] lie within the family of (M1), that is,
between a horizontal image and the container?
\item Do the product formula and the global realified Frobenioids do
more than leave one positive scalar common to all places?
\item Does a categorical statement, for instance
[IUTchIII, Cor.~2.3(i)] or [IUTchIII, Thm.~3.11(iii)], have a
consequence for regions or log-volumes that (M1) violates?
\end{enumerate}
\end{centralq}

The provisional answers in Appendix~\ref{s:questions} also remove the
difference of the theta and algebraic sections of [EtTh] as a
candidate (Q2, Q6).

In the decomposition of [Form, \S4], this asks where the arrow
(SHE), which the skeletal code takes as an input, is supplied.  In
the conditional assembly of [U26c], the carrier form of (S) is the
content of H*val; in the reading of [U26, Def.~4] it implies the
equality (9-1) of [LANA, \S9.2], which in turn implies the value
form [U26, Prop.~6].

Among the statements that the proof of Corollary~3.12 cites, none
passes (T1) in the reading of [U26]: the 22 that carry log-volume
content hold in its models [U26, Thms.~10, 22], and the others are
discussed in [U26, \S9].  This concerns the statements as they are
cited; what their proofs construct, in particular the proof of
Theorem~3.11, is what the central question asks about.
Other places in the source are listed in Appendix~\ref{s:places}.

\subsection{Forms of answer}\label{ss:answers}

Any of the following is useful, and one line is enough.
\begin{enumerate}[label=(\Alph*),leftmargin=*]
\item \emph{Misread.}  A reading above is wrong: one of (a)--(e)
does not hold in the models, because the actual construction is
larger there (where); or (M1) or (M2) violates a cited statement
(which one); or Step~(xi-e) does not use (SHE) in the form described
in Section~\ref{ss:step} (in which form, then).
\item \emph{Located.}  A property of the actual construction, not
among (a)--(e), that excludes (M1) or (M2), with a page, together
with one consequence for regions, log-volumes or labels that fails
in that model.
\item \emph{Not needed.}  The argument reaches (S) by a route other
than (xi-e)--(xi-f); please name the route.
\item \emph{Not yet written.}  The supplier is contained in the
construction but not yet stated in a citable form; please indicate
the part of the construction, for instance a part of the proof of
Theorem~3.11 or one of the triangles of [Form, \S4].
\end{enumerate}

\subsection{Strength is not evidence}\label{ss:strength}

\begin{remark}\label{r:strength}
A candidate supplier has, together with [IUTchIV, Thm.~1.10], the
strength of an explicit height inequality, because (S) has.  This
is a property of every proof of such an inequality: at the step
where its content enters, a proof carries the strength of its
conclusion.  The fact is therefore expected whether or not the
construction supplies the step.  It is not offered as evidence
either way, and it does not make the question hard to answer, since
the question asks where and not how.  It is used above only as a
criterion: a supplier, once found, can be recognized by what it
implies.  The component questions of Appendix~\ref{s:questions} are
weaker than the central question, and each informs a part of it.
\end{remark}

\section{Fixed background}\label{s:background}

\subsection{The residual in two-path form}

The residual in the conditional assembly [U26c] compares two
routes at the level of objects.  Route~1 starts from the
$\Theta$-pilot, passes through the value-group part of one
$\Theta_{\mathrm{LGP}}$-link to the $q$-pilot, and then uses the
column-one log-Kummer correspondence.  Route~2 starts from the
same $\Theta$-pilot, uses the column-zero log-Kummer
correspondence and permutation symmetry to obtain a possible
image.  The object-level assertion $(\alpha)$ says that the
Route~1 representation is, up to the permitted indeterminacies,
one of the Route~2 possible images.  The numerical assertion
$(\beta)$ follows from $(\alpha)$ together with the decomposition
identity expressing the $q$-side readout as the volume of that
object and monotonicity of the hull volume.  We use $(\alpha)$ as
the standard form of the residual.

\subsection{Parts treated as supplied}

The questions below fix the following layers as supplied at
assembly-statement level.  This is a focusing device, not a claim
that they have been proved unconditionally in the actual theory.

\begin{center}
\small
\begin{tabular}{@{}p{0.18\textwidth}p{0.55\textwidth}p{0.20\textwidth}@{}}
\toprule
Layer & Content & Status here\\
\midrule
Chart & Calibration, generator scale and transport distortion of the measuring chart & Treated as supplied\\
Unit & Theorem 3.11(iii)(a)(b), Kummer classes of actual functions and the restriction bridge & Treated as supplied\\
Semantics & Reading prime-strip morphisms as isomorphisms of collections of data & Fixed as a reading; actual inhabitation remains separate\\
Profile & Hull compatibility of the switch $\{j^2u\}_j$ to $\{u\}_j$ & \textbf{Open; subject of this note}\\
\bottomrule
\end{tabular}
\end{center}

\subsection{What the profile countermodel fixed}

The countermodel of [U26d] uses concrete data at $\ell=7$ in
which chart calibration, the normalized generator, the unit side,
membership and both volume decompositions hold, while only the
labelwise profile differs and the Corollary-3.12-shaped
conclusion fails.  The aggregate scale is in general
$12/(\ell(\ell+1))$.  Hence an input solving the residual cannot
be merely the numerical assumption that an aggregate scale equals
$1$; it must connect profiles at object level or at the level of
a sufficiently rich readout.

\subsection{Nearby public Lean work}

The public Lean development of Project LANA [lana] already contains
formalizations close in theme, including abstract forms of
[EtTh] Proposition~2.14, shift automorphisms, discrete and
cyclotomic rigidity, and, in Wall/LogLink models, scaling and
unit-induced properties, a profile coefficient \texttt{sFactor},
and hexagon noncommutativity.  This note does not restate those
developments.  Its questions concern how such model-level
classifications connect with the categories of the source texts
and with globalization of section differences.

\subsection{Formal cores checked in Lean, and what was not checked}\label{ss:cores}

Before circulation, the parts of the implication map that can be
stated without the [EtTh] or IUT constructions were isolated in
the Lean module:
\begin{center}\small
\texttt{M3.relative.ActualValuePartQuestionImplications}.
\end{center}
A selected build over three modules passed, and
\texttt{\#print axioms} for the main theorems used only the
standard three axioms.  The checked abstract implications are:
\begin{itemize}[leftmargin=*]
\item two splittings of one group homomorphism agree on the
quotient, their pointwise difference lies in the kernel and
satisfies the usual cocycle identity; hence any readout factoring
through that quotient agrees for the two splittings;
\item object membership, the readout decomposition identity, the
hull-domain condition and monotonicity imply the numerical volume
bound, abstractly $(\alpha)\Rightarrow(\beta)$;
\item equality of an otherwise unconstrained transported number
with the $q$-side readout does not imply that bound, and equal
totals do not determine a labelled profile;
\item a normalized linear readout invariant under all label
permutations is an average, the mean-square formula holds, and a
common ratio $r\ne1$ at every place does not cancel in a nonzero
total.
\end{itemize}

For this version, the elementary forms that the readings (a)--(e)
of Section~\ref{ss:test} use were checked in a standalone file that
depends only on Mathlib,
\begin{center}\small
\texttt{CentralQuestionReadings.lean},
\end{center}
compiled with Lean~4.30.0; for each of its 19 theorems,
\texttt{\#print axioms} shows only the standard three axioms.  It
checks the following.
\begin{itemize}[leftmargin=*]
\item Readings (a) and (b): an upper bound for the hull of the
maximal family is an upper bound for every smaller family; the hull
of a family that contains the container and lies in it is the hull
of the container; and an excess bounded independently of $N$ fails,
from some $N$ on, a requirement that grows linearly in $N$.
\item Readings (c) and (d): a statement about data that do not
contain the value of the $q$-pilot also holds in the configuration
that changes only that value, where (S) at the level of values
fails; so such a statement does not imply (S).
\item Reading (e), with Q0, Q1, Q2 and Q6: a scaling of the local
degrees that keeps the families of total degree zero is one scalar
common to all places; one scalar acting on all labels alike does not
carry the profile $\{a\}_j$ to $\{j^2a\}_j$, nor back, once two
labels are present; and neither does a transport that keeps every
degree readout.
\end{itemize}
Examples in the same file show that the positive slope, the
condition that the members lie in the container, the two labels and
the second place are needed, and removing any one of six hypotheses
makes the corresponding proof fail.  These are the forms that the
readings use.  That the constructions of [IUTchIII] and the models
of [U26] have these forms is the content of the readings, and it is
not formalized.

The Lean files do not formalize the readings of
Section~\ref{ss:test} themselves, mono-theta environments, the
actual difference $\delta$, its globalization, Statement~P$'$,
the two return bridges, or dictionary-dependent relations among
candidate forms.  They should not be read as verifying the
questions themselves.

\subsection{Relation to prior public claims, in the audited range}

The central identification difficulty itself is not new: it has
been a public issue since [SS], and Project LANA records a
corresponding membership problem as (9-1) [LANA, \S9.2].  Earlier
work on the same passage is listed in Section~\ref{ss:known}.  Joshi has published,
inside his own arithmetic Teichm\"uller theory, a claimed local
prototype of Corollary~3.12 and valuation-scaling results
[JosII].  Because the definitions and framework differ, this note
does not treat those claims, whose evaluation is not attempted
here, as directly closing Questions~Q1 or Q2 as phrased in the
[EtTh] categories.  For local outer automorphisms arising from
the Jannsen--Wingberg presentation, see [NSW], [HN] and [Kon];
in the statements audited here, the emphasis is on unit/log-side
actions, not a category-by-category classification of actions on
a normalized degree quotient or a one-place global realization.
A public document [Zen26] claims an additional hull-inflation
obstruction even under (9-1); that claim is not verified or used
here.  Accordingly, this note does not claim that Q0--Q7 are
globally unprecedented or open in every formulation.  In the
range audited, no numbered statement was located that closes
them in the exact categories and with the exact readouts used
below.

\section{Closing: what guidance would be decisive?}

The provisional answers in Appendix~\ref{s:questions} settle the
component questions as far as the cited statements allow.  Along the
route through the difference of sections, global synchronization
leaves only one positive scalar in the degree (Q0, Q1), the
difference of sections does not move degree readouts (Q2, Q5), and
the two-path difference is not that difference (Q6).

Any of the following would be decisive for this route:
\begin{itemize}[leftmargin=*]
\item for the central question, any answer of the forms (A)--(D) of Section~\ref{ss:answers}, in particular a pointer on one of the places (i)--(v) of Section~\ref{ss:test};
\item for the component questions, a correction of the provisional answers in Appendix~\ref{s:questions}, in particular those to Q0(4), Q1, Q2 and Q6, or a reference that settles them;
\item a prior public formulation of the residual in [EtTh] or anabelian reconstruction vocabulary.  This note does not make its value depend on a novelty claim.
\end{itemize}

The central question concerns a step that carries the strength
of the conclusion.  As Remark~\ref{r:strength} explains, this is
expected of the step at which the content of any proof of the
inequality enters, and it does not make the question hard to
answer.  The component questions are weaker, and each informs a part
of the central one.  The earlier works of Section~\ref{ss:known},
the conditional assembly and the profile countermodel have localized
the remaining difficulty to that step, and Section~\ref{ss:square}
locates it at one edge of the comparison: the value part of the
$\Theta$-link, carried to regions in the common container.  The purpose of this note is
to make the step searchable, and to put its parts in a form that
specialists can judge in everyday vocabulary.

\begin{center}
\small
\begin{tabular}{@{}p{0.12\textwidth}p{0.42\textwidth}p{0.39\textwidth}@{}}
\toprule
Residual & Content & Treatment here\\
\midrule
R1 & Numbered statements of IUT I--IV as black boxes & Not verified\\
R2V & Value-group profile switch H*val & Located at one edge (\S\ref{ss:square}) and asked in \S\ref{ss:test}; component questions answered provisionally (App.~\ref{s:questions})\\
R2 & Inhabitation / frame fidelity to the actual construction & Marked at Q6 and Form B seams\\
R3 & Source checking and independent re-verification & Author checking and outside rereading remain\\
\bottomrule
\end{tabular}
\end{center}

\appendix

\section{Setup and terminology}\label{s:setup}

Appendices~\ref{s:setup}--\ref{s:relations} keep the component
questions of earlier versions together with the setup that they use,
with provisional answers.  They are not needed to state or to answer
the central question of Section~\ref{s:central}.
Appendix~\ref{s:places} lists further places in the source.

\subsection{Stage}\label{ss:stage}

Let $K$ be a number field, $E/K$ an elliptic curve, and $S$ a
finite set of places at which $E$ has split multiplicative
reduction.  For $v\in S$, let $q_v\in K_v^\times$ be the Tate
parameter, with $\ord_v(q_v)>0$ for the normalized valuation.
Fix a prime $\ell\ge5$, a bound inherited as a hypothesis from
the assembly-side data conditions.  For each $v$, consider the
Tate curve, its $\ell$-torsion points, the relevant hyperbolic
fundamental groups, and decomposition groups at cusps and torsion
points.  Labels are $j\in\{1,\ldots,\ell^\divideontimes\}$,
$\ell^\divideontimes=(\ell-1)/2$.  Write $\eta_{v,j}$ for the
Kummer class of the theta value.  Its valuation satisfies
\[
\ord_v(\theta_{v,j})
 =\frac{j^2}{2\ell}\ord_v(q_v).
\]
With $u_v:=|\log(\uu q_v)|=(2\ell)^{-1}\log(\mathfrak q_v)$,
the signed readouts are $-j^2u_v$ on the theta side and $-u_v$ on
the $q$ side.

\subsection{Two sections and their difference}

A mono-theta environment in [EtTh] contains tempered fundamental
group data and a theta section.  The cyclotomic envelope also has
a tautological algebraic section.  In a relevant subquotient the
two appear as splittings of one surjection, with the theta
section obtained from the algebraic section by subtracting an
$\eta$-cocycle.  We write $\delta$ generically for such a
difference or for the change of section under a change of model,
always naming the category in question.  This avoids confusing
the standard $\eta$-cocycle of [EtTh] with indeterminacy under
conjugation or replacement of models.

A mono-theta environment contains the theta section as data; the
algebraic section is detected, in subquotients, through
commutators.  A bi-theta environment contains both sections but
loses discrete rigidity.  This tension is Question~Q3.

\subsection{What ``shear'' means here}

Because the word is used with different scopes, we use a working
definition.  For a replacement $\alpha$ of splittings, models or
identifications in a specified category, fixing specified
group-theoretic and cyclotomic data, consider the induced action
on degree readouts.  Conceptually there are two types:
\begin{center}
\small
\begin{tabular}{@{}p{0.20\textwidth}p{0.43\textwidth}p{0.30\textwidth}@{}}
\toprule
Type & Form of action & Effect on degree\\
\midrule
U (splitting shear) & Splitting changes by a unit-valued cocycle & Normally invisible in the valuation quotient\\
D (degree shear) & Normalized degree readout itself changes nontrivially & If linear, $D\mapsto\lambda D$, $\lambda\ne1$, is the issue\\
\bottomrule
\end{tabular}
\end{center}

A warning is essential.  Local automorphisms arising from the
Jannsen--Wingberg presentation and often called shears mix
generators and induce unipotent actions on the unit/log side
[NSW, Thm.~7.5.14], [HN], [Kon, Thm.~1.1].  Under suitable
hypotheses such an action need not be a scalar composed with a
field automorphism, and it may move sizes inside a log shell.
This should not be identified with a D-type action on the
normalized valuation quotient or degree readout itself.  Q0
therefore does not presuppose the existence of D-type shears; it
asks for their classification.

When an action on a generator readout is multiplicative, write
its ratio as $\lambda_v$.  Otherwise $\lambda_v$ is only shorthand
and the action itself should be recorded.  For transported
readouts write
$\varepsilon_v:=D_v^{\mathrm{tr}}(q_v)-D_v(q_v)$.

\section{The local skeleton already in [EtTh]}\label{s:skeleton}

This section records a starting point, not a new result.  The
theta functional equation [EtTh, Prop.~1.4(ii)] and its
cohomological form [EtTh, Prop.~1.5(iii)] are the source of the
$j^2$ profile.  By [EtTh, Thm.~1.6(ii)], a tempered isomorphism
preserves the kernel in the valuation chain and the element
$1\in\widehat{\mathbb Z}$; by (iii), the
$\mathcal O^\times\cdot\eta^\Theta$ class is preserved up to
conjugation.  Thus the indeterminacy of the theta class is
compressed to the unit side, while its value part is pinned.
Corollaries~2.8(iii) and 2.9 preserve the glued class without a
constant indeterminacy and preserve cusp labels.  Proposition~2.12
detects the algebraic splitting through commutators;
Proposition~2.14(ii) gives a shift automorphism inducing identity
on both quotient and kernel; and Corollary~2.19(i) obtains a
cyclotomic isomorphism from the difference of two splittings.

Formally, a readout factoring through the common quotient takes
the same value for the two splittings.  This abstract part is
among the Lean-checked cores.  It is only readout invariance, not
object identity, and it does not include identification of the
actual degree readout with such a quotient readout,
globalization, or transport to the procession.  The questions
below separate those outer steps.

\section{Component questions Q0--Q7}\label{s:questions}

Each of the following questions is weaker than the central
question of Section~\ref{s:central}, and none is needed to state
it.  They concern one route, through the difference of the theta and
algebraic sections of [EtTh], and they work at the level of the
degree at each place.  An action on the degree line of a place
changes the readouts of all labels by the same factor, so it cannot
by itself turn the profile $\{j^2u\}_j$ into $\{u\}_j$.  The component
questions can therefore close this route, or keep it open, but they
cannot by themselves supply (S); the profile enters at Q7, whose
distinction between readouts and membership is now part of (S).
Each question is followed by a provisional answer.  An answer marked
\emph{from the cited statements} rests on the statements it cites;
one marked \emph{author's reading} also rests on our reading of how
they apply, and corrections are welcome.  Taken together, the answers
close this route.
An answer to one of them would inform a specific part of the
central question, recorded under ``Bearing on the central
question''.  For each, a known reference, a countermodel, a
correction of category, or a better formulation would serve the
purpose of this note.  Structurally, Q0--Q5 cut out prerequisites,
definitions and the local core of Statement~P$'$ in
Appendix~\ref{s:pprime}.  Q6 and Q7 are bridges back from P$'$ to
the assembly.  Thus P$'$ is not an additional independent question
beside Q0--Q7; it is a composite form built from Q0--Q5, with
Q6--Q7 outside it.

We use the following labels for relations with (9-1):
\emph{directly connected} under a fixed dictionary; \emph{connected
under an additional identification}; \emph{requires separate
examination}; and \emph{corresponding}, meaning essentially the
same content in different vocabulary under a fixed dictionary.
Here ``dictionary'' means a fixed correspondence between the
assembly-side two paths, possible images and degree data and the
membership formulation used by LANA.  The dictionary itself is
not proved in this note.

\subsection*{Q0. In which category do local shears live?}

For a mono-theta environment over a local field and a degree
readout $D$, classify the actions on $D$ induced by automorphisms
or changes of model fixing specified group-theoretic and
cyclotomic data.  In particular, in which of the following
categories does a D-type action, moving degree itself, occur?
\begin{enumerate}[leftmargin=*]
\item topological-group automorphisms of tempered fundamental groups, modulo inner automorphisms;
\item automorphisms of mono-theta environments;
\item automorphisms of bi-theta environments;
\item automorphisms of tempered Frobenioids and their splitting monoids;
\item automorphisms of realified Frobenioids.
\end{enumerate}
Please also distinguish whether ``fixing cyclotomic data'' means
identity action on the cyclotomic object or compatibility with
its natural cyclotomic isomorphism.

\emph{Relation to (9-1): requires separate examination.}  (9-1)
concerns membership of specific standard objects and does not
classify automorphisms category by category.  Even if (9-1) is
true, degree-moving shears might exist in a wider category; even
if they exist, the standard transport might avoid them.

\emph{Bearing on the central question.}  A degree-moving action acts on
the degree line of a place, so it multiplies the readouts of all
labels by the same factor and cannot by itself turn $\{j^2u\}_j$ into
$\{u\}_j$.  Its bearing is on whether the indeterminacies contain a
freedom in the degree, that is, on how large the blurring is.  A
classification excluding such actions in a category removes that
category from consideration.

\emph{Answer (from the cited statements, except item~4).}  An
isomorphism of tempered fundamental groups induces isomorphisms of
the relevant cohomology groups that are compatible with the
surjections onto $\widehat{\mathbb Z}$ determined by the valuations,
preserving both their kernels and the element
$1\in\widehat{\mathbb Z}$ [EtTh, Thm.~1.6(ii)]; so no D-type action
occurs in category~1.  An automorphism of a mono-theta or a bi-theta
environment determines an automorphism of the tempered fundamental
group [EtTh, Prop.~2.14(iii)], so none occurs in categories~2 and~3
either.  These automorphisms act on the labels, through
$(l\cdot\mathbb Z)\rtimes\{\pm1\}$ in the mono-theta case and through
a subgroup of $(N^\dagger l\cdot\mathbb Z)\rtimes\{\pm1\}$ containing
$(Nl\cdot\mathbb Z)\rtimes\{\pm1\}$ in the bi-theta case.  After
realification the degree is determined only up to $\mathbb R_{>0}$
[FrdI, Thm.~6.4(ii)], so a D-type action occurs in category~5, by
one positive scalar; further structure [FrdI, Thm.~6.4(iii)(iv)]
reduces it to $\mathbb Q_{>0}$ or to $1$.  Category~4 is not settled
by the statements read here.  The two meanings of fixing cyclotomic
data give the same answers, since the degree is read through the
valuation surjection of [EtTh, Thm.~1.6(ii)].

\subsection*{Q1. Can a local shear globalize? (Main component question)}

Fix the local shear objects as in Q0.  What is the image, inside
the product of local shears over places, of differences arising
from a global change of model?  When the degree action can be
written by ratios $(\lambda_v)_{v\in S}$, can the family
\[
\lambda_{v_0}\ne1\quad\text{at one bad place \(v_0\) only},\qquad
\lambda_v=1\quad(v\ne v_0)
\]
be realized?  If not, is the obstruction only the product
formula, label synchronization, rigidity of a decomposition-group
preserving global model, or another global constraint?  We ask
for the distinction between a purely global anabelian reading,
using arithmetic fundamental-group data and decomposition groups,
and a synchronized reading in which label synchronization of the
type in [IUTchI, Prop.~6.5] is part of the given data.

If a one-place shear is realizable but compensated at other
places so that the total is preserved, the form of compensation
should be made explicit.  Formally, a common ratio $r\ne1$ at
all places does not cancel in a nonzero total; what remains open
is whether place-dependent families, and one-place support in
particular, occur in the image of global differences.

\emph{Relation to (9-1): requires separate examination.}  (9-1)
states membership for standard global objects; it does not
classify realizability of prescribed local families.  Under a
fixed dictionary and Bridge~(a), total consequences for the
standard difference may follow, but not the one-place support
condition.

\emph{Bearing on the central question.}  Tells whether a freedom in the
degree survives global synchronization.  Since such a freedom acts on
all labels alike, this bears on how large the indeterminacies are and
on the role of global data stressed in Step~(xii) of the proof of
Corollary~3.12 [IUTchIII, p.~185], rather than on a supplier
directly.

\emph{Answer (author's reading).}  By the answer to Q0, a
degree-moving action occurs only after realification, by a positive
scalar.  For the global realified Frobenioids of $F_{\mathrm{mod}}$,
compatibility with localization is stated up to $\mathbb R_{>0}$-orbits
[IUTchII, Cor.~4.10(v)], that is, up to one scalar for all places.  A
family with $\lambda_{v_0}\ne1$ and $\lambda_v=1$ for $v\ne v_0$ would
not preserve the product formula: it sends a principal divisor with a
nonzero component at $v_0$ to a divisor of nonzero total degree.  So
one-place support is not realizable, and the obstruction is the
product formula.  Such a scalar acts on all labels alike.

\subsection*{Q2. Is the degree readout invariant under the standard difference of sections?}

For the algebraic and theta sections in [EtTh], if the
identification of coefficients is transported to the theta-section
side, is the valuation-quotient readout of the $q$-parameter
Kummer class invariant?  Precisely, does
$\varepsilon_v=D_v^{\mathrm{tr}}(q_v)-D_v(q_v)$ vanish?  Is this
already a direct consequence of Propositions~2.12 and 2.14(ii),
Theorem~1.6(ii)(iii), and Corollary~2.19(i), or is an additional
claim needed at the coefficient-identification step?  Abstractly,
invariance is immediate once the readout factors through a
quotient on which the two splittings agree.  The remaining issue
is the source-level identification showing that the actual degree
readout and transport have that form.

\emph{Relation to (9-1): connected under an additional
identification.}  With the assembly dictionary and readout
identification fixed, a numerical shadow follows from (9-1); the
exact transformation rule inside [EtTh] does not.  This question
may close from existing [EtTh] results alone.

\emph{Bearing on the central question.}  With the dictionary and the
readout identification fixed, (9-1) implies a numerical shadow of
$\varepsilon_v=0$ (see the relation above), so a nonzero value would
close this route.  The value $0$ is consistent with the route, but it
does not supply the labelwise switch.

\emph{Answer (author's reading).}  $\varepsilon_v=0$, and it follows
from the statements cited in the question.  The difference of the two
splittings is taken on the subquotient over $l\cdot\Delta_\Theta$ and
yields the cyclotomic isomorphism [EtTh, Cor.~2.19(i)]; the
automorphism obtained by shifting by a $\boldsymbol\mu_N$-valued
difference of sections is the identity on both the quotient and the
kernel [EtTh, Prop.~2.14(ii)]; and the valuation readout of Kummer
classes of elements of the base field, such as $q_v$, is preserved by
tempered isomorphisms together with $1\in\widehat{\mathbb Z}$
[EtTh, Thm.~1.6(ii)].  Transporting the identification of
coefficients to the theta-section side therefore leaves the readout
of $q_v$ unchanged, and no further claim is needed at the
coefficient-identification step.  This agrees with the abstract
invariance checked in Lean (Section~\ref{ss:cores}).

\subsection*{Q3. What is the smallest category containing both sections and the degree readout?}

What is the smallest category in which the algebraic section,
the theta section and the degree readout simultaneously exist as
data while discrete rigidity is retained?  In a mono-theta
environment the algebraic section is only detected in
subquotients through commutators, whereas a bi-theta environment
has both sections but loses discrete rigidity.  Is there a
standard category resolving this tension, or evidence that the
tension itself is essential?

\emph{Relation to (9-1): requires separate examination.}  Even if
(9-1) is formulated in the LANA framework, that does not decide
which mono/bi-theta environment is the minimal home for this
question.  It is a question at the level of source definitions.

\emph{Bearing on the central question.}  Locates the smallest
category in which this route, and a supplier along it, could be
stated.

\emph{Answer (author's reading).}  The mono-theta environment
itself.  The algebraic section is not part of its data, but the
commutator property of [EtTh, Prop.~2.12] takes its place, and the
three rigidities of [EtTh, Cor.~2.19] hold; [EtTh, Rem.~2.19.2] notes
that the commutator property plays the role of the algebraic section,
which does not appear in a mono-theta environment.  Adding the
algebraic section as data, as in a bi-theta environment, cuts the
symmetry acting on the labels down to multiples of $N$
[EtTh, Prop.~2.14(iii)].  In this sense the tension is essential, and
[EtTh] resolves it by using the algebraic section as a property rather
than as data.  The degree readout is available through
[EtTh, Thm.~1.6(ii)].

\subsection*{Q4. Does a global $\delta$ exist without Hodge theatres?}

Can the projective system of local mono-theta environments and
its $\eta$-classes be glued, from global arithmetic
fundamental-group data and restrictions to decomposition groups
alone, into one global object $\delta$ whose local components
are not independent?  If so, what functor or reconstruction
algorithm gives the gluing?  If not, what is the smallest
additional datum required?  Is the understanding correct that
label synchronization alone does not give cocycle coherence?

\emph{Relation to (9-1): requires separate examination.}  A
global object in the LANA framework would still need to be
identified with a gluing of [EtTh] section systems.  Whether the
setup can be founded in pre-IUT vocabulary is not decided by
granting (9-1).

\emph{Bearing on the central question.}  Decides whether the global
datum that a supplier might use can be stated without Hodge
theatres.

\emph{Answer (from the cited statements).}  Not in the naive sense.
The source states that the local tempered fundamental groups at bad
places do not extend to a global tempered fundamental group
[IUTchIII, Cor.~3.12, proof, Step~(viii), p.~180].  What is used
instead is a global $\pm$-synchronization, which produces a profinite
conjugacy indeterminacy in the vertically coric construction of the
theta values at torsion points; this indeterminacy is resolved by the
theory of [IUTchI, \S2] ([IUTchI, Rem.~6.12.4(iii)];
[IUTchII, Rems.~4.5.3(iii), 4.11.2(iii)]).  So the additional datum
is this synchronization, which is part of the Hodge-theatre
apparatus.  We found no statement that makes the local
$\eta$-classes coherent as cocycles; the synchronization is at the
level of conjugacy.

\subsection*{Q5. Can the degree action of $\delta$ be computed in a Tate prototype?}

For a Tate curve over a finite $p$-adic field containing the
necessary roots of unity, for example at $\ell=5$ and including
Gaussian-type local test cases, make explicit from theta series
and product expansions the classes $\theta_{v,j}$ and the change
of $\eta$-classes under translation by $a\in\mathbb Z$.  Under
the standard splitting obtained by restriction to the zero
section, compute the transported degree action
$\varepsilon_v$, separating classically known parts from parts
depending on the choice of section.

\emph{Relation to (9-1): partly known; the centre requires
separate examination.}  Valuations of theta values and the
functional equation follow from classical computation and [EtTh,
Props.~1.4--1.5].  The value of $\varepsilon_v$ under the standard
splitting is not computed by (9-1).

\emph{Bearing on the central question.}  A computation that tests Q2
in a Tate prototype.

\emph{Answer (author's reading).}  The transported degree action is
$\varepsilon_v=0$, by the answer to Q2.  The profile itself comes from
the quadratic exponents of the theta series and its functional
equation
$\ddot\Theta(q_X^{a/2}\ddot U)=(-1)^aq_X^{-a^2/2}\ddot U^{-2a}\ddot\Theta(\ddot U)$
[EtTh, Prop.~1.4(ii)], through the fact that restricting the \'etale
theta class to a point gives the value of the theta function there up
to units [EtTh, Prop.~1.4(iii)]; this gives the valuations
$\frac{j^2}{2\ell}\ord_v(q_v)$ of Appendix~\ref{ss:stage}.  The
difference of sections does not contribute to the profile.  We did
not tabulate the values at $\ell=5$, since the answer does not depend
on them.

\subsection*{Q6. Bridge~(a): is the assembly two-path difference the [EtTh] section difference?}

At each bad place, restrict the composite
``Route~1 theta-link value part $\circ$ (Route~2 permutation
transport)${}^{-1}$'' to the data of the mono-theta splitting
monoid of [EtTh, Thm.~2.2(ii)].  Does it agree with transport
induced by the [EtTh] section difference, the $\eta$-cocycle?
Is that agreement compatible, including degree readouts, with
the passage from the single-basepoint method to the procession
method in [IUTchII, Rem.~3.5.3] and with horizontal compatibility
in [IUTchIII, Cor.~2.3(iv)]?  If this is immediate from a numbered
statement, we would be grateful for that reference.

\emph{Relation to (9-1): corresponding under a fixed dictionary.}
With a dictionary between LANA objects and the assembly two
paths, a positive answer largely overlaps the content of (9-1).
But fixing that dictionary is itself an unproved bridge, and
identification with the [EtTh] $\delta$ is not immediate from the
statement of (9-1).

\emph{Bearing on the central question.}  A positive answer would
reduce the central question, along this route, to properties of the
[EtTh] section difference; a negative answer would close the route.

\emph{Answer (author's reading; it depends on Q2).}  In the strict
sense, no.  The two-path difference sends the labelwise readouts
$\{j^2u_v\}_j$ to $\{u_v\}_j$; that is the profile switch.  The
transport induced by the [EtTh] section difference leaves the degree
readout of the $q$-datum unchanged (Q2).  The two can agree on degree
readouts only if all labels have the same readout, which fails for
$\ell\ge5$.  So this route does not carry the switch, and the
compatibility with [IUTchII, Rem.~3.5.3] and
[IUTchIII, Cor.~2.3(iv)] does not arise for it.  In the sense of
inclusion in the hull, agreement up to the indeterminacies is (S)
itself.

\subsection*{Q7. Bridge~(b): what is needed to pass from equal readouts to object membership?}

Under a splitting into unit part and a rational ray generated by
a normalized element, can two pilot objects with the same unit
part and the same total procession degree be treated as the same
possible image in the hull-compatible sense?  If not, what is
the smallest additional datum beyond total degree: the profile
tuple itself, a lower-bound or inclusion condition for regions,
or another object-level identification?

\emph{Relation to (9-1): membership is supplied; the implication
requires separate examination.}  Granting (9-1) supplies the
needed membership itself.  It does not prove an implication
``equal readouts $\Rightarrow$ membership''.  A total degree is a
many-to-one shadow of a tuple, and without additional structure
such an implication fails in general.  In the assembly-shaped
countermodel, this insufficiency is visible when the transported
number is not connected to a theta-side object.

\emph{Bearing on the central question.}  This is where the profile
enters.  The distinction between equal readouts and membership is
now part of (S) in Section~\ref{ss:target}: a supplier of the value
form need not supply the carrier form.

\emph{Answer (from the cited statements and Section~\ref{ss:cores}).}
No.  Equal unit parts and equal total procession degree do not make
two pilot objects the same possible image: a total is a many-to-one
shadow of a profile (Section~\ref{ss:cores}; the model of [U26d]).
For the value form of (S) an inequality of totals is enough; for
membership, that is, the carrier form of (S) and (9-1) in the reading
of [U26, Def.~4], an inclusion of regions is needed.  That inclusion
is the smallest additional datum, and it is what the central question
asks for.

\section{Statement \texorpdfstring{P$'$}{P'}}\label{s:pprime}

Statement P$'$ is a revised presentation of an earlier
reformulation, arranged to match the local skeleton and the
decomposition into questions above.  It is not claimed to be
equivalent to the residual.  P$'$ is not an additional
independent problem beside Q0--Q7: Q0--Q5 separate its
prerequisites, definitions and local core, while Q6--Q7 are
bridges outside P$'$ leading back to the assembly.

\begin{center}
\small
\begin{tabular}{@{}p{0.34\textwidth}p{0.20\textwidth}p{0.39\textwidth}@{}}
\toprule
Part of P$'$ & Questions & Remark\\
\midrule
Two sections and $\delta$ setup & Q3 & Minimal category\\
Replacement by degree readouts & Q0, Q2, Q5 & Shear classification, transported readout, computation\\
Globalization & Q1, Q4 & Image of global differences; existence of global $\delta$\\
Total compatibility assertion & Q1 standard case + Q2 & Standard difference, not arbitrary shear families\\
Return to assembly & Q6, Q7 & Bridges~(a)(b), outside P$'$\\
\bottomrule
\end{tabular}
\end{center}

\subsection{Differences from a solved form}

[EtTh, Cor.~2.19(i)] says that the difference of two sections
gives a compatible isomorphism on cyclotomic objects.  P$'$
differs in two ways: first, it replaces the object of
compatibility by value-group degree readouts; second, it states
the local theory globally, under one global $\delta$ and the
product-formula constraint.  The geometric Szpiro inequality in
the function-field case is only an expository anchor for the
second difference: there globality appears as a closed
intersection-theoretic total.  It does not provide a mechanism
in the number-field case.  For a prior formulation of a
geometric Corollary~3.12, see [Jos].

\subsection{Setup and statement}

Take $K$, $E/K$, $S$, $q_v$, $\ell\ge5$, theta torsion data and
$\eta_{v,j}$ as in Appendix~\ref{s:setup}.  Take the two sections
$s^{\mathrm{alg}}$, $s^\Theta$ and their difference $\delta$ in
the sense of Appendix~\ref{s:setup}; cyclotomic compatibility is known by
[EtTh, Cor.~2.19(i)].  As a marked
\textbf{premise-or-part-of-the-question}, suppose that the
sections and $\delta$ exist as one global object whose local
sections are restrictions to decomposition groups.  This
globalization is not founded here from the local [EtTh] theory;
it is Q4.

\begin{statement}[P$'$: compatibility of transported degree readouts]
$\delta$ is compatible with value-group degree readouts.
Precisely, for the system over all $v\in S$, the degree datum
$\ord_v(q_v)$ read through the algebraic section agrees, as the
same readout in the same reconstructed value group, with the
degree datum of the same object read through the theta section,
that is, from the system of Kummer classes $\{\eta_{v,j}\}$.
\end{statement}

What is read through the theta section is not the total degree
of the theta system itself, which is computed from its own
valuations and does not equal the $q$-degree.  It is the degree
of $q$-data transported by $\delta$.  Nor is the assertion about
capsulewise profiles $\{j^2\}$ versus $\{1\}$, or a place-by-place
assertion.  It is asserted only for the system over $S$, for one
global $\delta$, under the product-formula constraint.

\emph{Relation to (9-1): connected under additional
identifications.}  With dictionary and Bridges~(a)(b) fixed, the
numerical content of P$'$ is obtained as a shadow of (9-1).  The
converse does not hold: P$'$ is a readout assertion, not object
membership itself.

\subsection{Rating of \texorpdfstring{P$'$}{P'}}

No equivalence between P$'$ and $(\alpha)$ has been proved; the
gap consists of Q6 and Q7.  The globalization in the setup is an
unverified seam, and a version with the global clause removed is
not P$'$.  The local skeleton gives formal scaffolding, but
transplantation to degree readouts is itself the open content
and does not follow automatically from Corollary~2.19(i), whose
proof uses cyclotomic Kummer structure.  The direction checked
formally is $(\alpha)$ plus decomposition identities implying
$(\beta)$; numerical compatibility P$'$ alone does not imply
$(\alpha)$.  If P$'$ is false, this assembly route closes, but
that would not disprove ABC.  If P$'$ is true, Bridges~(a)(b)
and the IUT black box still remain.  With the provisional answers of
Appendix~\ref{s:questions}, the local part of P$'$ holds (Q2), but
Bridge~(a) fails in the strict sense (Q6), so P$'$ does not lead to
$(\alpha)$ along this route.

\section{Candidate forms A--D and the implication map}\label{s:forms}

Before this question list, the residual $(\alpha)$ was examined
in the following candidate forms.  We record only their relation
to the questions here.  The implications are sketches under a
fixed dictionary, not proved equivalences between autonomous
statements.

\begin{center}
\small
\begin{tabular}{@{}p{0.10\textwidth}p{0.34\textwidth}p{0.49\textwidth}@{}}
\toprule
Form & Content & Relation here\\
\midrule
A & Commutativity, in the value component, of the two-path square & Assembly-side form of $(\alpha)$; Q6 asks for [EtTh]-side identification\\
B & Application of semantic principle S to the relevant isomorphism and value component & Semantic clothing of A; the general principle is not the question here\\
C & Commutativity of mono-anabelian reconstruction $R$ with theta transport & Algorithmic clothing of A; seam: assembly degree data are output of $R$\\
D $=$ (9-1) & Membership form of $\eta^q=\eta^{\mathrm{anab}}_S$ type & Current reference form; this note asks before and around it\\
P$'$ & Global compatibility of transported degree readouts & Common-vocabulary entrance; weaker readout form than A\\
\bottomrule
\end{tabular}
\end{center}

Under supplied inputs, the sketch is: the relevant instance of B
corresponds to A and to $(\alpha)$; C would give A through its
seam; and D corresponds to $(\alpha)$ under the fixed dictionary.
These dictionary-dependent relations were not formalized in Lean.
Q0--Q5 separate prerequisites and definitions outside this map;
Q6--Q7 ask about its seams.

\section{Summary of relations with (9-1)}\label{s:relations}

The labels below mean the position of a question if (9-1) is
granted.  They do not include the truth of (9-1) or correctness
of a dictionary.  This is a division-of-roles summary.  Many
questions retain independent meaning after (9-1), as category
classification or foundations for globalization.

\begin{center}
\small
\begin{tabular}{@{}p{0.30\textwidth}p{0.25\textwidth}p{0.38\textwidth}@{}}
\toprule
Question & Relation & Reason in brief\\
\midrule
Central question (\S\ref{s:central}) & Supplier of (9-1) & A supplier of the carrier form of (S) gives (9-1) in the reading of [U26, Def.~4]; (9-1) gives the value form\\
Q0 local-shear classification & Separate & (9-1) gives no categorywise automorphism classification\\
Q1 globalization / one place & Separate & Standard membership does not classify prescribed families\\
Q2 transported-readout invariance & Additional identification & Numerical shadow follows; internal [EtTh] rule is separate\\
Q3 minimal category & Separate & Source-definition question\\
Q4 global $\delta$ & Separate & Identification with [EtTh] gluing is additional\\
Q5 Tate computation & Partial & Theta valuations known; $\varepsilon_v$ is not computed by (9-1)\\
Q6 Bridge~(a) & Corresponding under dictionary & Dictionary itself unproved\\
Q7 Bridge~(b) & Membership supplied; implication separate & (9-1) does not prove readout-to-membership\\
P$'$ & Additional identifications & Shadow under dictionary and bridges; converse fails\\
Forms A, D & Corresponding under dictionary & Rephrasings of $(\alpha)$\\
Form B, general principle & Separate & One instance does not imply a general semantic principle\\
Form C & Additional identification & Assembly degree data as output of $R$ is the seam\\
\bottomrule
\end{tabular}
\end{center}

\section{Other places in the source}\label{s:places}

Besides the constructions named in the central question, the
following places may bear on the edge of Section~\ref{ss:square}.
The list is not exhaustive.

{\small
\begin{longtable}{@{}>{\raggedright\arraybackslash}p{0.27\textwidth}>{\raggedright\arraybackslash}p{0.37\textwidth}>{\raggedright\arraybackslash}p{0.30\textwidth}@{}}
\toprule
Place & Content & Bearing\\
\midrule
\endhead
{}[IUTchIII, Thm.~3.11(i)], pp.~153--155 & Multiradial representation; at the end of (i), the permutation symmetries of the \'etale-like structures & How the output is made independent of the holomorphic structure\\\addlinespace[3pt]
{}[IUTchIII, Thm.~3.11(ii)], pp.~155--156 & Log-Kummer correspondence; (Ind3) & Step (xi-d) calls the column-one log-Kummer correspondence essential\\\addlinespace[3pt]
{}[IUTchIII, Thm.~3.11(iii)], pp.~156--158 & Compatibility of the Kummer isomorphisms with the $\Theta^{\times\boldsymbol\mu}_{\mathrm{LGP}}$-links & Relates the $\Theta$-link to Kummer theory\\\addlinespace[3pt]
{}[IUTchIII, Rem.~3.11.1(iii)--(v)], pp.~160--163 & (IPL), (SHE), (APT) and the hidden internal structure & The properties that Step~(xi) keeps, and what it sets aside\\\addlinespace[3pt]
{}[IUTchIII, Rem.~3.11.1(vi)--(viii)], pp.~163--167 & Naive models: distinct labels give (IPL) but leave (SHE) unclear; confusing $q$ with its powers gives a trivial multiradial representation; the features of Theorem~3.11 on which the estimate depends; the permutation symmetry $\Pi\to G\leftarrow\Pi$ & The source's own statement of the dilemma, and of how the theory is meant to avoid it\\\addlinespace[3pt]
{}[IUTchIII, Rem.~3.9.5(vii)--(x)], pp.~131--145 & Hull, (Ob1)--(Ob9), (sQ4), the closed loop of (viii)(ix) & The passage from regions to comparable numbers, cited in (xi-c)--(xi-e)\\\addlinespace[3pt]
{}[IUTchIII, Cor.~3.12, proof], Steps (i)--(x), pp.~175--181, and (xii), p.~185 & Set-up of the comparison; in (xii), the role of the global realified Frobenioids in the $\Theta^{\times\boldsymbol\mu}_{\mathrm{LGP}}$-link & Step (xii) stresses that global data matter in Step (xi)\\\addlinespace[3pt]
{}[IUTchIII, Rem.~3.12.2], pp.~187--195 & Interpretation of the gluing and of the entanglement of addition and multiplication & Context of (xi-a)\\\addlinespace[3pt]
{}[IUTchII, Cor.~4.5(iv)(v)], p.~133; [IUTchII, Cor.~4.6(iv)(v)], pp.~138--139; [IUTchII, Cor.~4.10(v)], p.~160 & Theta and Gaussian monoids and Frobenioids, the ratio vector $(j^2)_j$, compatibility up to $\mathbb R_{>0}$-orbits & The source of the profile\\\addlinespace[3pt]
{}[EtTh, Thm.~1.6, Props.~2.12, 2.14, Cor.~2.19] & Mono-theta environments, the difference of the theta and algebraic sections, cyclotomic rigidity & The local skeleton of Appendix~\ref{s:skeleton}; part of the structure set aside in (xi-b)\\\addlinespace[3pt]
{}[FrdI, Thm.~6.4] & Restoration of the degree scale in tiers: $\mathbb R_{>0}$, $\mathbb Q_{>0}$, $1$ & What an equivalence of realified Frobenioids fixes: a single scalar on one line\\\addlinespace[3pt]
{}[Form, \S\S3--4] & Symmetry of the \'etale Hodge theatre under exchanging the domain and the codomain of the $\Theta$-link with the basic prime-strip fixed; the four triangles; (APT) in the second and third & The author's own decomposition of the step\\
\bottomrule
\end{longtable}}

\section*{Disclosure}

This is personal work by an independent researcher, unrelated to
any institutional affiliation.  Generative AI was used under the
author's direction for auditing, formalization and drafting.
Responsibility for the contents rests with the author.  This is
a work report not yet independently re-verified.

\begin{thebibliography}{99}\small
\bibitem[IUTchI]{iut1} S.~Mochizuki, Inter-universal Teichm\"uller Theory I: Construction of Hodge Theaters, \emph{Publ. Res. Inst. Math. Sci.} 57 (2021), 3--207.
\bibitem[IUTchII]{iut2} S.~Mochizuki, Inter-universal Teichm\"uller Theory II: Hodge-Arakelov-theoretic Evaluation, \emph{Publ. Res. Inst. Math. Sci.} 57 (2021), 209--401.
\bibitem[IUTchIII]{iut3} S.~Mochizuki, Inter-universal Teichm\"uller Theory III: Canonical Splittings of the Log-theta-lattice, \emph{Publ. Res. Inst. Math. Sci.} 57 (2021), 403--626.
\bibitem[IUTchIV]{iut4} S.~Mochizuki, Inter-universal Teichm\"uller Theory IV: Log-volume Computations and Set-theoretic Foundations, \emph{Publ. Res. Inst. Math. Sci.} 57 (2021), 627--723.
\bibitem[EtTh]{etth} S.~Mochizuki, The \'Etale Theta Function and its Frobenioid-theoretic Manifestations, \emph{Publ. Res. Inst. Math. Sci.} 45 (2009), 227--349.
\bibitem[FrdI]{frd1} S.~Mochizuki, The Geometry of Frobenioids I: The General Theory, \emph{Kyushu J. Math.} 62 (2008), 293--400.
\bibitem[FrdII]{frd2} S.~Mochizuki, The Geometry of Frobenioids II: Poly-Frobenioids, \emph{Kyushu J. Math.} 63 (2009), 1--182.
\bibitem[AbsTopIII]{abstop3} S.~Mochizuki, Topics in Absolute Anabelian Geometry III: Global Reconstruction Algorithms, \emph{J. Math. Sci. Univ. Tokyo} 22 (2015), 939--1156.
\bibitem[AbsAnab]{absanab} S.~Mochizuki, The Absolute Anabelian Geometry of Hyperbolic Curves, in \emph{Galois Groups and Fundamental Groups}, MSRI Publ.~41, Cambridge Univ. Press, 2003, 77--122.
\bibitem[SS]{ss} P.~Scholze and J.~Stix, Why abc is still a conjecture, 2018.
\bibitem[NSW]{nsw} J.~Neukirch, A.~Schmidt, and K.~Wingberg, \emph{Cohomology of Number Fields}, 2nd ed., Springer, 2008, Thm.~7.5.14.
\bibitem[HN]{hn} Y.~Hoshi and Y.~Nishio, On the outer automorphism groups of the absolute Galois groups of mixed-characteristic local fields, \emph{Res. Number Theory} 8 (2022), no.~3, Paper No.~56.
\bibitem[Kon]{kon} K.~Kondo, Anabelian aspects of the outer automorphism groups of the absolute Galois groups of mixed-characteristic local fields, arXiv:2512.09231v2.
\bibitem[Jos]{jos} K.~Joshi, Construction of Arithmetic Teichmuller Spaces II$_{1/2}$: Deformations of Number Fields, arXiv:2305.10398.
\bibitem[JosII]{jos2} K.~Joshi, Construction of Arithmetic Teichmuller Spaces II: Proof of a local prototype of Mochizuki's Corollary 3.12, arXiv:2303.01662v3, Thms.~6.9.1 and 9.2.1. Not treated here as evaluated input.
\bibitem[Zen26]{zen} C.~Zeng, An Integral-Structure Inflation in the Hull Operation of IUT Corollary 3.12, Hexagon 2610.00646v1. Not verified or used as input here.
\bibitem[LANA]{lana} Project LANA, \emph{Project LANA interim report on IUT theory}, July 2026, 50 pp. Mentioned as a reference form; its results are not used here as proved inputs.
\bibitem[lana]{lanarepo} Project LANA, \emph{lana-agents/iut}, Lean repository, \url{https://github.com/lana-agents/iut}; files \texttt{Iut/WallEquiv/NonDerivable.lean} (added 13 October 2026, commit 180c049) and \texttt{Iut/Wall/Constraints.lean}, and the EtaleTheta, Wall and LogLink modules, consulted at commit 553076a (3 November 2026). Mentioned as nearby formalization; its results are not used here as proved inputs.
\bibitem[Higuchi]{higuchi} J.~R.~Higuchi, A formal negative result for the passage from Step (xi-e) to Step (xi-f) in IUT III, preprint, March 31, 2026, 9 pp., DOI \href{https://doi.org/10.13140/RG.2.2.10989.76006}{10.13140/RG.2.2.10989.76006}.
\bibitem[Pajo]{pajo} P.~Pajo, Scale-free common-container comparisons do not entail numerical inequalities: a minimum-viable no-go theorem motivated by the IUT III Theorem 3.11-to-Corollary 3.12 bridge, preliminary research note, April 26, 2026, 14 pp., DOI \href{https://doi.org/10.13140/RG.2.2.33620.16005}{10.13140/RG.2.2.33620.16005}.
\bibitem[Takkun]{takkun} Takkun-kohinata, \emph{IUT\_LEAN}, Lean repository, \url{https://github.com/Takkun-kohinata/IUT_LEAN}, commit 9d56c46 (15 July 2026); files \texttt{IutLean/Multiradial/ThetaHullDichotomy.lean} and \texttt{DisputedCore.lean}.
\bibitem[Saito]{saito} Y.~Saito, A registered-provenance audit of IUT III, Corollary 3.12, Step (xi), version 2.0.0, Zenodo, 26 July 2026, DOI \href{https://doi.org/10.5281/zenodo.21596198}{10.5281/zenodo.21596198}.
\bibitem[sfingali]{sfingali} sfingali, The Scholze--Stix objection, formally: the arithmetic is agreed, the $q$-side is absent, Zenodo, 2 August 2026, DOI \href{https://doi.org/10.5281/zenodo.21757640}{10.5281/zenodo.21757640}.
\bibitem[Geiger]{geiger} L.~Geiger, The bridge certificate: a line-level audit instrument for inter-theoretic transfer arguments, Zenodo, first version 14 August 2026, DOI \href{https://doi.org/10.5281/zenodo.21925803}{10.5281/zenodo.21925803} (all versions).
\bibitem[Form]{form} S.~Mochizuki, On the formalization of IUT: a preliminary progress report (joint work in progress with Y.~Hoshi, G.~Yamashita, Y.~Yang), slides, April 2026, 14 pp.
\bibitem[Ri]{ri} D.~Richards, \emph{motizukipdf --- Mochizuki's papers, plus a toolkit for the data types of IUT}, GitHub repository \url{https://github.com/oneman/motizukipdf}; Lean port \texttt{lean/IUT/LeanForm.lean} at commit 2665db1 (30 October 2026).
\bibitem[DHPS]{dhps} T.~Dupuy and A.~Hilado, Probabilistic Szpiro, baby Szpiro, and explicit Szpiro from Mochizuki's Corollary 3.12, arXiv:2004.13108v2 [math.NT], 29 April 2020.
\bibitem[U26]{u26} T.~Urahigashi, Where Theorem 3.11 sits, Zenodo, DOI \href{https://doi.org/10.5281/zenodo.24104932}{10.5281/zenodo.24104932} (v1.2.1).
\bibitem[U26b]{u26b} T.~Urahigashi, Forgetting and Completeness, Zenodo, DOI \href{https://doi.org/10.5281/zenodo.24210290}{10.5281/zenodo.24210290} (v1.1.5).
\bibitem[U26c]{u26c} T.~Urahigashi, A conditional assembly from Theorem 3.11 to Corollary 3.12: decomposition of the scalar layer, a semantic principle, and isolation of the value-group profile switch, Zenodo, DOI \href{https://doi.org/10.5281/zenodo.24209900}{10.5281/zenodo.24209900} (v1.0.3).
\bibitem[U26d]{u26d} T.~Urahigashi, Scale preservation is not enough: a profile-level independence countermodel for the Corollary 3.12 assembly inputs, Zenodo, DOI \href{https://doi.org/10.5281/zenodo.24209904}{10.5281/zenodo.24209904} (v1.0.1).
\bibitem[SL]{sl} T.~Urahigashi, Segment or loop? A homotopy-type-theoretic overview of the dispute over Corollary 3.12 of inter-universal Teichm\"uller theory, version 1.4.0, Zenodo, DOI \href{https://doi.org/10.5281/zenodo.24039360}{10.5281/zenodo.24039360}.
\end{thebibliography}

\end{document}




